% Convergence analysis and finite-trajectory projection bounds.
% Detailed proofs: supplement/S4_theory_details.tex.
\ifdefined\CGAChapterFiveEmbedded
\else
\documentclass[11pt]{article}
\usepackage[a4paper,margin=25mm]{geometry}
\usepackage{amsmath,amssymb,amsthm,mathtools,booktabs,microtype,enumitem}
\usepackage[hidelinks]{hyperref}
\usepackage{cleveref}
\def\CGAChapterFiveStandalone{1}
\newtheorem{theorem}{Theorem}[section]
\newtheorem{lemma}[theorem]{Lemma}
\newtheorem{proposition}[theorem]{Proposition}
\newtheorem{corollary}[theorem]{Corollary}
\newtheorem{assumption}[theorem]{Assumption}
\theoremstyle{remark}\newtheorem{remark}[theorem]{Remark}
\begin{document}\setcounter{section}{3}
\fi
\providecommand{\cE}{\mathcal E}
\providecommand{\cD}{\mathcal D}
\providecommand{\cP}{\mathcal P}
\providecommand{\cK}{\mathcal K}
\providecommand{\cnorm}[1]{\left\lVert#1\right\rVert}
\providecommand{\cabs}[1]{\left|#1\right|}
\providecommand{\abs}[1]{\left|#1\right|}
\providecommand{\cinner}[2]{\left\langle#1,#2\right\rangle}
\providecommand{\cR}{\mathbb R}
\providecommand{\cdd}{\,\mathrm d}
\ifdefined\CGAChapterFiveStandalone
  \newcommand{\CGAHLDreference}{the finite-pool descent estimates}
  \newcommand{\CGAEnergyMetricReference}{the energy--distance theorem and metric corollary in the main text}
  \newcommand{\CGAEnergyTheoremReference}{the energy--distance theorem in the main text}
\else
  \newcommand{\CGAHLDreference}{Section~\ref{sec:c5-descent}}
  \newcommand{\CGAEnergyMetricReference}{Theorem~\ref{thm:energy-v} and Corollary~\ref{cor:metric-rates}}
  \newcommand{\CGAEnergyTheoremReference}{Theorem~\ref{thm:energy-v}}
\fi

\section{Convergence analysis and finite-trajectory projection bounds}
\label{sec:c5}

We first derive finite-pool descent estimates in terms of pool coverage, score evaluation, radial stationarity, and coefficient correction. Local Hessian transfer then connects nonlinear selection with approximation in a frozen Hilbert metric. Finally, source representations and normalized innovations give projection bounds for nested spaces, which the quartic energy identity converts into nonlinear energy estimates.

Throughout, $m$ counts accepted independent directions, $m_0$ is a starting
index, $\widehat G_m$ is a computed state, $G_m^\sharp$ an exact active-space
minimizer, and $z_m$ a frozen orthogonal projection.  All derivatives act on
the specified feasible linear space.

\subsection{Finite-pool descent and the atomic energy rate}
\label{sec:c5-descent}
\label{sec:descent}

Let \(X\) be a real Banach space and \(\cE:X\to(-\infty,+\infty]\) a proper convex energy with minimizer \(u^*\). Assume that \(\cE\) is finite and continuously differentiable on an open neighborhood containing the states and dictionary segments used below. For a symmetric unit dictionary \(\cD_X\), set
\[
 \widehat\delta_m=\cE(\widehat G_m)-\cE(u^*),\qquad
 S_m^X=\sup_{g\in\cD_X}|\cinner{\cE'(\widehat G_m)}g|.
\]
For a finite pool $\cP_m\subset\cD_X$ and $0<t_m^{\rm pool}\le1$, write
\[
 |\widetilde s_m(g)-s_m(g)|\le\xi_m,\qquad
 \widetilde s_m(g_{m+1})\ge t_m^{\rm pool}\max_{g\in\cP_m}\widetilde s_m(g),
\]
where the score $s_m$ and the coverage factor $\tau_m$ are defined in
\Cref{alg:finite} and \eqref{eq:pool-coverage}, respectively.
\begin{assumption}[Directional quadratic upper bound]
\label[assumption]{ass:directional}
For some $L>0$, every $m$, every $g\in\cD_X$, and $|\lambda|\leq1$,
\begin{equation}
 \cE(\widehat G_m+\lambda g)\leq\cE(\widehat G_m)+\lambda\cinner{\cE'(\widehat G_m)}g+\tfrac L2\lambda^2,
\label{eq:c5-directional}
\end{equation}
\end{assumption}
Assume that the expanded active space contains the signed step used below and that
the inexact active correction satisfies
\[
 \cE(\widehat G_{m+1})\leq
 \inf_{v\in\widehat\Phi_{m+1}}\cE(v)+\varepsilon_m^{\rm opt},
 \qquad \varepsilon_m^{\rm opt}\geq0.
\]
The radial active-space stationarity defect is denoted by
\begin{equation}
 \abs{\cinner{\cE'(\widehat G_m)}{\widehat G_m}}
 \leq \eta_m^{\rm stat}.
 \label{eq:c5-orthogonality-defect}
\end{equation}
This is the radial component of the active-space stationarity residual; the
full active-space dual KKT residual is a separate quantity.

\begin{proposition}[Finite-pool descent]
\label{prop:c5-descent}
\label{thm:inexact}
Under these conditions,
\begin{equation}
 \widehat\delta_{m+1}\leq\widehat\delta_m-H_L(a_m)+\varepsilon_m^{\rm opt},\qquad
 a_m=\left[t_m^{\rm pool}\tau_m S_m^X-(1+t_m^{\rm pool})\xi_m\right]_+,
 \label{eq:c5-descent}
\end{equation}
where
\[
  H_L(a)=\max_{0\le\lambda\le1}
    \left(\lambda a-\frac{L}{2}\lambda^2\right),
  \qquad a\ge0.
\]
\end{proposition}
\begin{proof}
The score error is inserted before and after the finite-pool maximization.  Testing the enlarged active space with a signed step of length $\min\{1,a_m/L\}$ gives the displayed decrease, and the correction error is added once.
\end{proof}

If $B_{\rm src}>0$, $u^*/B_{\rm src}$ belongs to the closed absolutely convex hull of $\cD_X$, and the
radial defect satisfies \eqref{eq:c5-orthogonality-defect}, convexity gives
\begin{equation}
 \widehat\delta_m\leq\eta_m^{\rm stat}+B_{\rm src}S_m^X.
 \label{eq:c5-atomic}
\end{equation}
This relation bounds the energy defect by the dictionary score and the radial stationarity residual.

\begin{lemma}[Nonautonomous scalar recurrence]
\label{lem:c5-recurrence}
If $x_m\ge0$ and $x_{m+1}\le x_m-b_mx_m^q$ with $q>1$ and $b_m\ge0$, then, for
$x_0>0$,
\begin{equation}
 x_n\le\left[x_0^{1-q}+(q-1)\sum_{m=0}^{n-1}b_m\right]^{-1/(q-1)}.
 \label{eq:c5-recurrence}
\end{equation}
If $x_0=0$, the sequence vanishes.
\end{lemma}
\begin{proof}
Apply convexity to $x\mapsto x^{1-q}$ while the sequence is positive and sum the resulting increments; once a term vanishes, nonnegativity gives the remaining cases.
\end{proof}

\begin{corollary}[Atomic energy rate]
\label{cor:c5-atomic-rate}
Assume the hypotheses of \Cref{prop:c5-descent} and
$u^*\in B_{\rm src}K_1(\cD_X)$ with $B_{\rm src}>0$.
Suppose that, for some $m_0$, $\underline t>0$, and
$0\le\vartheta,\zeta<1$, the radial defect satisfies
\eqref{eq:c5-orthogonality-defect} and
\[
  t_m^{\rm pool}\tau_m\ge\underline t,\qquad
  \eta_m^{\rm stat}
    +\frac{B_{\rm src}(1+t_m^{\rm pool})}
                 {t_m^{\rm pool}\tau_m}\xi_m
    \le\vartheta\widehat\delta_m,\qquad
  \varepsilon_m^{\rm opt}\le\zeta H_L(a_m)
\]
for all $m\ge m_0$.
If $\widehat\delta_{m_0}>0$, set
\[
  b=\frac{\underline t(1-\vartheta)}{B_{\rm src}},
  \qquad A=b\widehat\delta_{m_0},
  \qquad c=\frac{(1-\zeta)b^2}{2\max\{L,A\}}.
\]
Then
\[
  \widehat\delta_{m_0+j}
    \le\left(\widehat\delta_{m_0}^{-1}+cj\right)^{-1},
  \qquad j\ge0,
\]
and hence
$\widehat\delta_m=O((m-m_0+1)^{-1})$.
If $\widehat\delta_{m_0}=0$, then
$\widehat\delta_m=0$ for every $m\ge m_0$.
\end{corollary}

\begin{proof}
Proposition~\ref{prop:c5-descent} and the correction-error bound give
\[
  \widehat\delta_{m+1}
    \le\widehat\delta_m-(1-\zeta)H_L(a_m),
\]
so the energy errors are nonincreasing.
Equation~\eqref{eq:c5-atomic} and the relative-error condition imply
$a_m\ge b\widehat\delta_m$.
Since $0\le b\widehat\delta_m\le A$ and $H_L$ is nondecreasing,
\[
  H_L(a_m)\ge H_L(b\widehat\delta_m)
    \ge\frac{b^2\widehat\delta_m^2}{2\max\{L,A\}}.
\]
Thus
$\widehat\delta_{m+1}\le
  \widehat\delta_m-c\widehat\delta_m^2$.
Apply Lemma~\ref{lem:c5-recurrence} to
$x_j=\widehat\delta_{m_0+j}$ with $q=2$.
The zero initial-error case follows from monotonicity.
\end{proof}

The ideal continuous method is recovered by setting
$\tau_m=1$ and
$\xi_m=\eta_m^{\rm stat}=\varepsilon_m^{\rm opt}=0$,
with selection strength bounded below.

For the standard pure, regularized, and reaction $p$-energies,
\CGAEnergyMetricReference{} give energy--distance equivalence
and the $W^{1,p}$ conversion. Hence an energy estimate
$\widehat\delta_m=O(m^{-\alpha})$ implies
\[
  d_{V,\bullet}(\widehat G_m,u^*)=O(m^{-\alpha/2}),\qquad
  \|\widehat G_m-u^*\|_{W^{1,p}}=O(m^{-\alpha/p}).
\]

\subsection{Local Hessian transfer}
\label{sec:c5-hessian}

The local transfer result connects nonlinear energy minimization with Hilbert-space approximation. Let $H$ be a
Hilbert space on which $\cE$ is twice Fr\'echet differentiable in an open
convex neighborhood of its minimizer $u^*$. Let $A_*=\cE''(u^*)$ be bounded,
symmetric, and coercive, and set
$a_*(v,w)=\langle A_*v,w\rangle$, $\|v\|_*^2=a_*(v,v)$.
Assume the closed ball $\|u-u^*\|_*\le R_*$ lies in this neighborhood and
\begin{equation}
 |\langle[\cE''(u)-A_*]v,w\rangle|
 \le L_*\|u-u^*\|_*\|v\|_*\|w\|_*,\qquad L_*R_*\le\tfrac12 .
 \label{eq:c5-hessian}
\end{equation}
For a finite-dimensional subspace $\Phi\subset H$, let
$P_\Phi:H\to\Phi$ be the $a_*$-orthogonal projection, characterized by
\[
  a_*(u-P_\Phi u,v)=0
  \qquad (u\in H,\ v\in\Phi).
\]
Set $z_\Phi=P_\Phi u^*$, and let $G_\Phi^\sharp$ be the exact
minimizer of $\cE$ over $\Phi$.
Assume that both $z_\Phi$ and $G_\Phi^\sharp$ lie in this ball.

\begin{lemma}[Taylor and projection estimates]
\label{lem:c5-projection}
Under \eqref{eq:c5-hessian}, for $u$ in this ball,
\begin{align}
 \left|\cE(u)-\cE(u^*)-\tfrac12\|u-u^*\|_*^2\right|&\le\tfrac{L_*}{6}\|u-u^*\|_*^3,\\
 \|\cE'(u)-A_*(u-u^*)\|_{*'}&\le\tfrac{L_*}{2}\|u-u^*\|_*^2,\\
 \|G_\Phi^\sharp-z_\Phi\|_*&\le L_*\|z_\Phi-u^*\|_*^2.
 \label{eq:c5-proj}
\end{align}
\end{lemma}
\begin{proof}
Taylor's formula gives the first two inequalities.  The third follows by testing the active-space first-order condition with $G_\Phi^\sharp-z_\Phi$ and absorbing the term $L_*R_*\|G_\Phi^\sharp-z_\Phi\|_*^2$.
\end{proof}

Assume the algorithmic dictionary $\cD_X\subset H$ satisfies
\begin{equation}
 0<c_*\le\|g\|_*\le C_*<\infty\quad(g\in\cD_X),
\qquad \cD_* = \left\{g/\|g\|_*:g\in\cD_X\right\}.
 \label{eq:c5-bridge}
\end{equation}
Let $\Phi_{m+1}=\Phi_m+\operatorname{span}\{g_{m+1}\}$, where
$g_{m+1}\in\cD_X\setminus\Phi_m$, and set $P_m=P_{\Phi_m}$.
Set $z_m=P_mu^*$ and define
$S_m^*=\sup_{d\in\cD_*}|\langle A_*(z_m-u^*),d\rangle|$.  At the exact
active-space minimizer $G_m^\sharp$ define the score defect
\[
 \eta_m^H=\sup_{g\in\cD_X}
 \left|\langle\cE'(G_m^\sharp),g\rangle-a_*(z_m-u^*,g)\right|.
\]
Lemma~\ref{lem:c5-projection} yields
$\eta_m^H\le C_{T}\|z_m-u^*\|_*^2$ on the stated ball, with $C_T$ uniform
over the specified dictionary.

\begin{proposition}[Relative score transfer]
\label{prop:c5-transfer}
If nonlinear selection at $G_m^\sharp$ over $\cD_X$ has strength
$t\in(0,1]$ and $0<t_0<tc_*/C_*$, then
\begin{equation}
 (1+t)\eta_m^H\le (tc_*-t_0C_*)S_m^*
 \label{eq:c5-relative}
\end{equation}
implies that the selected direction satisfies
\[
 \left|a_*(z_m-u^*,g_{m+1}/\|g_{m+1}\|_*)\right|\ge t_0S_m^*.
\]
\end{proposition}
\begin{proof}
The selected unnormalized frozen score is bounded below by
$t c_*S_m^*-(1+t)\eta_m^H$; division by the selected frozen norm costs at most
$C_*$.  This proves the claim.
\end{proof}

To make the entropy-to-rate step explicit, fix a restart index $m_0$ and set
$Q_0=I-P_{m_0}$, $H_0=Q_0H$, $f_0=Q_0u^*$, and
$\cD_{\rm rst}=Q_0\cD_*\setminus\{0\}$.  The projected atoms retain their projected norms.
Suppose $K=K_1(\cD_{\rm rst})$ is compact and $f_0\in B_*K$, $B_*>0$.
For weak OGA of strength $t_0$ in $H_0$, let
$d_i\in\cD_{\rm rst}$ be the selected atoms and
$\Pi_j$ the orthogonal projection onto the span of the first
$j$ selected atoms, with $\Pi_0=0$.
The bounded-atom entropy argument of
\cite[Theorem~4.2]{LiSiegelEntropy2024} gives
\begin{equation}
 \|f_0-\Pi_jf_0\|_*
 \leq \frac{B_*}{t_0}\frac{(j!V_j)^{1/j}}{\sqrt j}
 \varepsilon_{j+1}(K;H_0),
 \qquad V_j=\frac{\pi^{j/2}}{\Gamma(j/2+1)}.
 \label{eq:c5-entropy}
\end{equation}
Indeed, writing $y_i=\|f_0-\Pi_if_0\|_*^2$ and
$\ell_i^{\rm rst}=\|(I-\Pi_{i-1})d_i\|_*$, the source and weak selection give
\[
 y_i\le y_{i-1}-\frac{t_0^2y_{i-1}^2}{B_*^2(\ell_i^{\rm rst})^2},
 \qquad
 y_j\le\frac{B_*^2}{t_0^2j}
       \left(\prod_{i=1}^j\ell_i^{\rm rst}\right)^{2/j}.
\]
The second bound follows by reciprocation and the arithmetic--geometric mean
inequality while the errors are positive. The cross-polytope generated by
$\{\pm d_1,\ldots,\pm d_j\}$ lies in $\Pi_jK$ and has volume
$(2^j/j!)\prod_i\ell_i^{\rm rst}$.
A cover of $K$ by $2^j$ balls yields
$\prod_i\ell_i^{\rm rst}\le j!V_j\varepsilon_{j+1}(K;H_0)^j$,
by taking the infimum over covering radii. This proves \eqref{eq:c5-entropy};
zero errors persist. Stirling's formula bounds $(j!V_j)^{1/j}/\sqrt j$
uniformly.

\begin{theorem}[Hessian transfer after restart]
\label{thm:c5-hessian}
Assume \eqref{eq:c5-hessian}, \eqref{eq:c5-bridge}, and exact active-space
minimization. At each step $m_0\le m<m_0+J$, suppose nonlinear selection has
strength at least $t$ and \eqref{eq:c5-relative} holds with the same $t_0>0$.
Assume the exact minimizers and frozen projections remain in the Hessian ball
for $m_0\le m\le m_0+J$, $K$ is compact, and
\[
 f_0\in B_*K_1(\cD_{\rm rst}),\qquad
 \varepsilon_j(K_1(\cD_{\rm rst});H_0)\le C_{\rm ent}j^{-\beta},
 \quad 1\le j\le J+1,
\]
with $B_*,C_{\rm ent}>0$, $\beta>0$. Then, for $1\le j\le J$,
\begin{align}
 \|z_{m_0+j}-u^*\|_*&\le C B_*C_{\rm ent}t_0^{-1}j^{-\beta},\\
 \cE(G_{m_0+j}^\sharp)-\cE(u^*)&\le C_{\rm nl}j^{-2\beta}.
 \label{eq:c5-hessian-rate}
\end{align}
Here \(C\) is absolute and \(C_{\rm nl}\) depends on \(L_*,R_*,B_*,C_{\rm ent},t_0\).
\end{theorem}
\begin{proof}
The relative score proposition gives weak frozen selection after projection by
$Q_0$, so the restarted spaces form weak OGA on $H_0$.  Apply
\eqref{eq:c5-entropy}.  The projection lemma gives
$\|G_{m_0+j}^\sharp-u^*\|_*\le(1+L_*R_*)\|z_{m_0+j}-u^*\|_*$;
the Taylor estimate then yields the energy bound.
\end{proof}

\subsection{Approximate-source and quartic consequences}
\label{sec:c5-source}

For the finite-trajectory analysis, let \(Y\) be a finite-dimensional real Hilbert space containing \(u^*\) and the active spaces in the prescribed range. Let \(a_*(\cdot,\cdot)\) be its frozen inner product and \(P_m\) the associated orthogonal projection onto \(\widehat\Phi_m\). Assume
$\widehat\Phi_{m+1}=\widehat\Phi_m+\operatorname{span}\{d_{m+1}\}$ with
$d_{m+1}\notin\widehat\Phi_m$, and set
\[
 z_m=P_mu^*,\qquad e_m=u^*-z_m,\qquad r_m=\|e_m\|_*.
\]
Let $\cD_{\rm fin}\subset Y$ be a finite symmetric dictionary containing every accepted direction in this range.
For the selected direction $d_{m+1}$ define its innovation norm
$\ell_m=\|(I-P_m)d_{m+1}\|_*>0$ and
$S_m^{\rm fin}=\max_{d\in\cD_{\rm fin}}|a_*(e_m,d)|$.
Set $\theta_m=|a_*(e_m,d_{m+1})|/S_m^{\rm fin}$ when
$S_m^{\rm fin}>0$, and $\theta_m=0$ otherwise.  Assume that the restarted
residual has the explicit approximate atomic representation
\[
 Q_0u^*=\sum_{d\in\cD_{\rm fin}}a_dQ_0d+\rho_0,
 \qquad \sum_d|a_d|\leq B_{\rm entry},\qquad \|\rho_0\|_*\leq\sigma,
 \quad Q_0=I-P_{m_0}.
\]
For any $0\leq\underline\theta_m\leq\theta_m$, define
$W_j=\sum_{m=m_0}^{m_0+j-1}\underline\theta_m^2/\ell_m^2$.

\begin{theorem}[Weighted finite-trajectory source bound]
\label{thm:c5-source}
For $x_0=(r_{m_0}^2-\sigma^2)_+>0$ and $B_{\rm entry}>0$,
\begin{equation}
 (r_{m_0+j}^2-\sigma^2)_+
 \leq\left(x_0^{-1}+\frac{W_j}{4B_{\rm entry}^2}\right)^{-1},
 \qquad 0\le j\le J.
 \label{eq:c5-source-initial}
\end{equation}
If $x_0=0$ or $B_{\rm entry}=0$, monotonicity gives $r_m\leq\sigma$ on the range;
for $W_j=0$ the estimate is interpreted through projection monotonicity.
\end{theorem}
\begin{proof}
Pythagoras gives \(r_m^2-r_{m+1}^2=|a_*(e_m,d_{m+1})|^2/\ell_m^2\), and the source representation gives \(S_m^{\rm fin}\ge(r_m^2-\sigma^2)_+/(2B_{\rm entry})\). When both \(r_m>\sigma\) and \(r_{m+1}>\sigma\), set \(x_m=r_m^2-\sigma^2\). Then
\[
 x_{m+1}\le x_m-\frac{\underline\theta_m^2x_m^2}{4B_{\rm entry}^2\ell_m^2}.
\]
Taking reciprocals and summing gives the bound up to the first crossing of \(\sigma\). At and after that crossing, projection monotonicity gives \((r_m^2-\sigma^2)_+=0\).
\end{proof}

As a concrete example of converting projection bounds into nonlinear
energy estimates, consider the one-dimensional pure and regularized
quartic models. Let
\[
 X=\{u\in W^{1,4}(0,1):\textstyle\int_0^1u=0\},\qquad
 \cE_\varepsilon(u)=\int_0^1
 \left(\tfrac{\varepsilon^2}{2}|u'|^2+\tfrac14|u'|^4\right)\cdd x-\ell(u),
\]
and assume $u^*\in X$ solves its weak equation.  Let $Y\subset X$ contain
these finite-dimensional spaces and define
\[
 a_*(v,w)=\int_0^1(\varepsilon^2+3|(u^*)'|^2)v'w'\cdd x.
\]
For $\varepsilon>0$ this is positive definite on $Y$; for $\varepsilon=0$ assume
it is positive definite on all of $Y$.  In this paragraph
$\widehat\delta_m=\cE_\varepsilon(\widehat G_m)-\cE_\varepsilon(u^*)$.  Then
\begin{equation}
 \widehat\delta_m=\tfrac12(r_m^2+D_m^2)+R_m,
 \qquad
 R_m=\int_0^1\left((u^*)'v_m^3+\tfrac14v_m^4\right)\cdd x,
 \label{eq:c5-quartic}
\end{equation}
where $D_m=\|\widehat G_m-z_m\|_*$ and $v_m=\widehat G_m'-(u^*)'$. The identity yields the following two-sided energy comparison.
\begin{proposition}[Quartic energy comparison]
\label{prop:c5-quartic}
If $q_{\rm corr}\ge0$, $\rho\ge0$, $D_m\le q_{\rm corr}r_m$, and
$|R_m|\le\rho(r_m^2+D_m^2)$ with $\rho<1/2$, then
\begin{equation}
 (\tfrac12-\rho)r_m^2\le\widehat\delta_m\le
 (\tfrac12+\rho)(1+q_{\rm corr}^2)r_m^2.
 \label{eq:c5-quartic-comparison}
\end{equation}
\end{proposition}
\begin{proof}
Expand the quartic density about $(u^*)'$, cancel the linear term with the weak equation, and use orthogonality of the frozen projection to obtain \eqref{eq:c5-quartic}; the comparison follows immediately.
\end{proof}

\begin{theorem}[Power-type finite-range bound]
\label{thm:c5-power}
Define the normalized frozen innovation decrease
\[
 \gamma_m=\frac{|a_*(e_m,d_{m+1})|}{\ell_m},\qquad
 r_{m+1}^2=r_m^2-\gamma_m^2.
\]
Let $\alpha>0$, $c_0>0$, and suppose $r_{m_0}>0$.  If
$\gamma_m^2\ge c_0r_m^{2+2/\alpha}$ on
$m_0\le m<m_0+J$, then
\begin{equation}
 r_{m_0+j}^2\le\left(r_{m_0}^{-2/\alpha}+\frac{c_0j}{\alpha}\right)^{-\alpha},
 \qquad 0\le j\le J.
 \label{eq:c5-power}
\end{equation}
\end{theorem}
\begin{proof}
For positive successive residuals, apply convexity to
\(x^{-1/\alpha}\), with \(x_m=r_m^2\), and sum.
Once a residual is zero, nestedness keeps the subsequent residuals zero.
\end{proof}

\begin{corollary}[Energy and metric bounds for quartic models]
\label{cor:c5-power-metrics}
Under the assumptions of \Cref{thm:c5-power}, consider the
one-dimensional pure or regularized quartic energy defined above.
Suppose that the hypotheses of \Cref{prop:c5-quartic} hold for every
$m_0\le m\le m_0+J$, with common constants
$q_{\rm corr}\ge0$ and $0\le\rho<1/2$.
Let $U_j$ denote the right-hand side of \eqref{eq:c5-power}, and set
\[
  C_E=(1/2+\rho)(1+q_{\rm corr}^2).
\]
Then, for $0\le j\le J$,
\[
  \widehat\delta_{m_0+j}\le C_EU_j,\qquad
  d_{V,\bullet}(\widehat G_{m_0+j},u^*)
    \le C_V(C_EU_j)^{1/2},
\]
\[
  \|\widehat G_{m_0+j}-u^*\|_{W^{1,4}}
    \le C_S(C_EU_j)^{1/4}.
\]
Here $d_{V,\bullet}$ is the natural distance for the chosen model,
and $C_V,C_S>0$ are independent of $j$ on the stated range.
\end{corollary}

\begin{proof}
Combine \eqref{eq:c5-power} with
\eqref{eq:c5-quartic-comparison}, and apply
\CGAEnergyTheoremReference{}.
\end{proof}

\ifdefined\CGAChapterFiveEmbedded
\else
\end{document}
\fi
